\documentclass[11pt]{article}
\usepackage[margin=1in]{geometry}
\usepackage{amsmath,amssymb,amsthm}
\usepackage[colorlinks=true,urlcolor=blue]{hyperref}
\newtheorem{theorem}{Theorem}
\newcommand{\F}{\mathbb F}
\newcommand{\one}{\mathbf1}
\newcommand{\rank}{\operatorname{rank}}
\title{Degree masks do not restore a polynomial\protect\\intertwiner moment for nonregular lifts}
\author{Independent degree-refinement research}
\date{November 15, 2026}
\begin{document}
\maketitle

\paragraph{Scope.}
The regular-frame recovery packet uses a polynomial bound on the number
of finite-field intertwiners and enumerates them before checking the
integer one-hot equation. We prove that the corresponding moment is
exponential for an elementary nonregular base, even after all observable
degree/color masks and the linear one-hot row-sum constraints are added.
This is a deterministic obstruction to that algorithm, not by itself an
identifiability counterexample or a lower bound against other algorithms.

Let $H=K_{a,b}$, with fixed integers $b>a\ge3$. Name its part of $b$
vertices $L$ and its part of $a$ vertices $R$. Vertices in $L$ have degree
$a$, while those in $R$ have degree $b$. Let $G$ be \emph{any} $N$-lift
of $H$, with planted indicator $J$, and order its vertices and the base
vertices by these two parts. Write
\[
 A=\begin{pmatrix}0&B\\B^T&0\end{pmatrix},\qquad
 Q=\begin{pmatrix}0&\one_b\one_a^T\\
                       \one_a\one_b^T&0\end{pmatrix},
\]
where $B$ has size $bN$ by $aN$, and $AJ=JQ$.
Every row of $B$ sums to $a$ and every column to $b$.

\begin{theorem}\label{thm:obstruction}
Over any finite field $\F$ of size $p$, the affine set of matrices $Y$
satisfying all three requirements
\[
 AY=YQ,\qquad Y\one=\one,\qquad
 Y_{v,j}=0\ \text{if }\deg_G(v)\ne\deg_H(j)
\]
has cardinality at least
\[
 p^{(b-a)(b-1)N}.
\]
If the characteristic does not divide $a$, the same bound holds after
also requiring the planted column masses modulo that characteristic:
\[
 \one^TY=\one^TJ.
\]
For characteristic dividing $a$, that last additional constraint still
leaves at least $p^{((b-a)N-1)(b-1)}$ solutions when the exponent is
nonnegative.
\end{theorem}
\begin{proof}
Put
\[
 U=\ker B^T\subseteq\F^{bN},\qquad
 Z=\{z\in\F^b:\one_b^Tz=0\}.
\]
Rank-nullity gives $\dim U\ge(b-a)N$, and $\dim Z=b-1$ over every
field. For any $x\in U,z\in Z$, place $xz^T$ in the $L$-row,
$L$-column block of a matrix $V$ and put zero in the other blocks.
Then $AV=0$ because $B^Tx=0$, and $VQ=0$ because $z^T\one_b=0$.
Also $V\one=0$, and every nonzero entry respects the degree mask.
The span of these matrices is the tensor-product space $U\otimes Z$,
of dimension $(\dim U)(b-1)$. Adding any such matrix to $J$ gives an
admissible $Y$, proving the first bound.

If the characteristic does not divide $a$, every $x\in U$ has
\[
 a\one_{bN}^Tx=(B\one_{aN})^Tx
              =\one_{aN}^TB^Tx=0,
\]
so $\one^TV=0$ automatically. In the remaining characteristic,
restrict $U$ to $U\cap\one^\perp$, losing at most one dimension.
The same tensor argument gives the final bound.
\end{proof}

\paragraph{Degree refinement and observable colors.}
Starting from an uncolored graph, one round of color refinement
distinguishes the degree-$a$ and degree-$b$ parts. At every subsequent
round every vertex of the first part has exactly $a$ neighbors of the
second color, and every vertex of the second has exactly $b$ neighbors
of the first. Induction shows that color refinement has exactly these
two classes, for every lift and every $N$. The displayed degree mask
therefore already incorporates all colors obtainable in this way.
Refining by the multiset of neighbor degrees or any further ordinary
color-refinement rounds adds no constraint.

\paragraph{The smallest unbalanced example.}
Take $a=3,b=4$ and $\F=\F_2$. On every $N$-lift of $K_{3,4}$, even
after the degree/color mask, affine row sums and column-mass parity
constraints, there are at least $2^{3N}$ solutions. Since $G$ has
$7N$ vertices, this is $2^{3|V(G)|/7}$. Enumerating the entire affine
solution space and only then checking exact integer one-hot rows has
exponential running time on \emph{every} input in this lift model.
Its expected running time cannot be polynomial. In particular, the
regular packet's polynomial intertwiner moment cannot transfer to all
nonregular frames by simply observing vertex degrees first.

The matrices counted here need not have one-hot rows over the integers.
An algorithm which exploits the nonlinear one-hot constraints without
enumerating this affine space is not ruled out. Nor does this argument
show that $K_{3,4}$ fibers fail to be identifiable. Information-theoretic
obstructions for other fixed nonregular bases require a separate proof.
No historical firstness claim is made.
\end{document}
